\documentclass[11pt]{article}

\usepackage[margin=1in]{geometry}
\usepackage{amsmath,amssymb,amsthm}
\usepackage{bm}
\usepackage{enumitem}
\usepackage{booktabs}
\usepackage{graphicx}
\usepackage[colorlinks=true,linkcolor=blue,citecolor=blue]{hyperref}

%% Short-hands
\newcommand{\vv}[1]{\bm{#1}}          % bold vector / tensor
\newcommand{\nn}{\vv{n}}               % unit surface normal
\newcommand{\NN}{\vv{N}}               % stress-resultant tensor (bold)
\newcommand{\Nb}{\mathbf{N}}           % same, upright bold (for components)
\newcommand{\kone}{\kappa_1}
\newcommand{\ktwo}{\kappa_2}
\newcommand{\Hmean}{H}
\newcommand{\Kgauss}{K}
\newcommand{\divs}{\nabla_{\!s}\!\cdot}
\newcommand{\tp}{\otimes}              % tensor product
\newcommand{\R}{\mathbb{R}}

\title{Tension inference on closed surface meshes:\\
\large surface geometry, membrane equilibrium, and the GFDM solve}
\author{Working notes (vedo / spatchcocking / GFDM pipeline)}
\date{2026-06-17}

\begin{document}
\maketitle
\tableofcontents
\bigskip

%% =========================================================================
\section{Surface geometry and local frame}
\label{sec:geometry}

%% -------------------------------------------------------------------------
\subsection{Parametrisation and covariant basis}

Let $\Gamma\subset\R^3$ be a smooth, orientable surface (the midsurface of the
membrane wall).  We parametrise $\Gamma$ by the mapping
\begin{equation}
  \vv{x} = \varphi(\bm{\xi}), \qquad \bm{\xi}=(\xi^1,\xi^2)\in\Gamma_0\subset\R^2.
  \label{eq:param}
\end{equation}
The \emph{covariant basis vectors} tangent to $\Gamma$ are
\begin{equation}
  \vv{e}_a = \frac{\partial\vv{x}}{\partial\xi^a}, \qquad a\in\{1,2\}.
  \label{eq:cov-basis}
\end{equation}
These span the tangent plane $T_{\vv x}\Gamma$ at every point.  The
\emph{first fundamental form} (metric tensor) is
\begin{equation}
  g_{ab} = \vv{e}_a\cdot\vv{e}_b.
  \label{eq:metric}
\end{equation}
The unit outward normal is
\begin{equation}
  \nn = \frac{\vv{e}_1\times\vv{e}_2}{|\vv{e}_1\times\vv{e}_2|}.
  \label{eq:normal}
\end{equation}

%% -------------------------------------------------------------------------
\subsection{Second fundamental form and shape operator}

The \emph{second fundamental form} encodes how the normal $\nn$ tilts as one
moves along the surface:
\begin{equation}
  b_{ab} = \nn\cdot\frac{\partial^2\vv{x}}{\partial\xi^a\,\partial\xi^b}
          = -\frac{\partial\nn}{\partial\xi^a}\cdot\vv{e}_b.
  \label{eq:IIF}
\end{equation}
(Both expressions are equivalent by the Weingarten equation.)
The \emph{shape operator} (Weingarten map) is the mixed tensor
$W^a{}_b = g^{ac}b_{cb}$, whose eigenvalues are the \emph{principal curvatures}
$\kone\ge\ktwo$ and whose eigenvectors define the \emph{principal directions}.

The mean and Gaussian curvatures are
\begin{equation}
  \Hmean = \tfrac12(\kone+\ktwo), \qquad
  \Kgauss = \kone\ktwo.
  \label{eq:HK}
\end{equation}

%% -------------------------------------------------------------------------
\subsection{Principal curvature frame}
\label{sec:frame}

At a non-umbilic point ($\kone\ne\ktwo$) there exists a unique pair of mutually
orthogonal unit tangent vectors $\vv{e}_1^{\rm pc},\vv{e}_2^{\rm pc}$ aligned
with the principal directions.  Together with the outward normal, they form the
\emph{principal curvature frame}:
\begin{equation}
  \boxed{%
  \bigl(\,\vv{e}_1^{\rm pc},\;\vv{e}_2^{\rm pc},\;\nn\,\bigr),
  \qquad
  \vv{e}_1^{\rm pc}\times\vv{e}_2^{\rm pc}=\nn,
  }
  \label{eq:pc-frame}
\end{equation}
an orthonormal right-handed basis of $\R^3$.  At an umbilic point
($\kone=\ktwo$) any orthonormal pair in the tangent plane serves as the frame.

%% =========================================================================
\section{Membrane stress model}
\label{sec:model}

%% -------------------------------------------------------------------------
\subsection{Assumptions}

\begin{enumerate}[label=(\roman*)]
  \item \textbf{Thin wall.}  Wall thickness $t$ satisfies
        $t\ll\min(R_1,R_2)$ with $R_i=1/|\kappa_i|$, so bending moments are
        negligible and stress is uniform through the thickness.
  \item \textbf{Membrane state of stress.}  The wall stress is purely in-plane
        (tangential to $\Gamma$); there are no transverse shear or normal
        stresses.
  \item \textbf{Quasi-static equilibrium.}  Inertial forces are zero.
  \item \textbf{No tangential body force.}  The only applied load is a
        normal pressure jump $\Delta p = p_{\rm in}-p_{\rm out}>0$.
\end{enumerate}

%% -------------------------------------------------------------------------
\subsection{Stress resultant tensor}

The Cauchy stress integrated through the thickness gives the \emph{membrane
stress-resultant tensor} (force per unit length):
\begin{equation}
  \Nb = \int_{-t/2}^{+t/2}\bm{\sigma}\,ds,
  \qquad \Nb\,\nn = \vv{0}.
  \label{eq:resultant}
\end{equation}
The constraint $\Nb\,\nn=\vv{0}$ states that $\Nb$ is tangential.  In the
principal curvature frame it takes the form
\begin{equation}
  \Nb
    = N_{11}\,\vv{e}_1^{\rm pc}\tp\vv{e}_1^{\rm pc}
    + N_{22}\,\vv{e}_2^{\rm pc}\tp\vv{e}_2^{\rm pc}
    + N_{12}\bigl(\vv{e}_1^{\rm pc}\tp\vv{e}_2^{\rm pc}
                 +\vv{e}_2^{\rm pc}\tp\vv{e}_1^{\rm pc}\bigr),
  \label{eq:N-decomp}
\end{equation}
with three independent scalar DOFs $(N_{11},N_{22},N_{12})$ per vertex.  The
corresponding \emph{principal membrane stresses} are
\begin{equation}
  \sigma_1 = \frac{N_{11}}{t}, \qquad \sigma_2 = \frac{N_{22}}{t}.
  \label{eq:sigma}
\end{equation}

%% =========================================================================
\section{Balance of linear momentum}
\label{sec:balance}

%% -------------------------------------------------------------------------
\subsection{Vector equilibrium equation}

The membrane equilibrium follows from the balance of forces on an infinitesimal
surface patch.  In Cauchy form on the surface:
\begin{equation}
  \divs\Nb + \Delta p\,\nn = \vv{0}.
  \label{eq:equil}
\end{equation}
Here $\divs$ is the \emph{surface divergence}: for a surface tensor field
$\Nb$, $(\divs\Nb)^i = \nabla_\alpha N^{\alpha i}$ where $\nabla_\alpha$ is
the covariant surface derivative along $\vv{e}_\alpha$.

Equation~\eqref{eq:equil} is a vector equation in $\R^3$.  Projecting onto
$\nn$ and onto the two tangent directions separates it into one \emph{normal}
and two \emph{tangential} scalar equations.

%% -------------------------------------------------------------------------
\subsection{Normal component --- the Laplace (Young--Laplace) law}
\label{sec:normal}

Taking the dot product of~\eqref{eq:equil} with $\nn$ and using the identity
$(\divs\Nb)\cdot\nn = -\Nb:B$ (where $B=[b_{ab}]$ is the second fundamental
form regarded as a bilinear form), we obtain
\begin{equation}
  \Nb : B = \Delta p,
  \label{eq:laplace-tensor}
\end{equation}
i.e.\ $N^{ab}b_{ab}=\Delta p$.  In the principal curvature frame
($b_{ab}=\mathrm{diag}(\kone,\ktwo)$ and $N_{12}=0$ on the principal axes) this
reduces to the classical \textbf{Young--Laplace law}:
\begin{equation}
  \boxed{N_{11}\,\kone + N_{22}\,\ktwo = \Delta p,}
  \label{eq:laplace}
\end{equation}
or equivalently, in terms of principal stresses and thickness,
\begin{equation}
  \sigma_1\,\kone + \sigma_2\,\ktwo = \frac{\Delta p}{t}.
  \label{eq:laplace-sigma}
\end{equation}
This is the \emph{normal projection} of the equilibrium: one scalar equation
per surface point relating the two unknown stress components to the pressure
load and the local curvatures.

%% -------------------------------------------------------------------------
\subsection{Tangential components --- in-plane equilibrium}
\label{sec:tangential}

Projecting~\eqref{eq:equil} onto the two tangent directions $\vv{e}_\alpha^{\rm
pc}$ ($\alpha=1,2$) and using $\Delta p\,\nn\cdot\vv{e}_\alpha^{\rm pc}=0$ gives
\begin{equation}
  \bigl(\divs\Nb\bigr)\cdot\vv{e}_\alpha^{\rm pc} = 0, \qquad \alpha=1,2.
  \label{eq:tang-equil}
\end{equation}
In local coordinates these expand to:
\begin{align}
  \frac{\partial N_{11}}{\partial s_1}
  + \frac{\partial N_{12}}{\partial s_2}
  + (\kone-\ktwo)\,N_{12}\,c_{\alpha\beta}
  &= 0, \label{eq:tang1} \\[4pt]
  \frac{\partial N_{12}}{\partial s_1}
  + \frac{\partial N_{22}}{\partial s_2}
  + (\kone-\ktwo)\,N_{12}\,c_{\alpha\beta}
  &= 0, \label{eq:tang2}
\end{align}
where $\partial/\partial s_\alpha$ denotes the arc-length derivative along
$\vv{e}_\alpha^{\rm pc}$ and the curvature-coupling terms arise from the
Christoffel symbols of the surface connection.

\paragraph{On closed surfaces.}  Equations~\eqref{eq:laplace} and
\eqref{eq:tang-equil} together form a $3n$-dimensional system for $(N_{11},
N_{22},N_{12})$ at every vertex.  The system is \emph{statically indeterminate}
on a closed surface (genus 0): there exists a finite-dimensional space of
self-equilibrating (null) stress states.  These are suppressed by Tikhonov
regularisation (Section~\ref{sec:gfdm}).

%% =========================================================================
\section{Ambient-component form and the GFDM discretisation}
\label{sec:gfdm}

%% -------------------------------------------------------------------------
\subsection{Working in world Cartesian components}

A key simplification arises when $\Nb$ is expressed in its \emph{ambient
Cartesian components} $N^{ij}$ ($i,j\in\{1,2,3\}$) rather than covariant
surface components.  The surface divergence then reads
\begin{equation}
  \bigl(\divs\Nb\bigr)^i
    = \sum_j\Bigl[
        (v_1)_j\,\partial_\xi N^{ij} + (v_2)_j\,\partial_\eta N^{ij}
      \Bigr],
  \label{eq:divs-ambient}
\end{equation}
where $(v_1,v_2,\nn)$ is any local orthonormal tangent frame and
$\partial_\xi,\partial_\eta$ are the surface derivatives along $v_1,v_2$.  The
Christoffel symbols (surface connection terms) are \emph{absorbed into the
ambient representation}: they do not appear explicitly, which makes the
formulation valid on arbitrary surfaces without coordinate charts.

\paragraph{Constraint $\Nb\,\nn=\vv 0$.}  In the local frame $(v_1,v_2)$ the
constraint is automatically satisfied by the parameterisation
\begin{equation}
  \Nb = p\,(v_1\tp v_1) + q\,(v_2\tp v_2) + r\,(v_1\tp v_2 + v_2\tp v_1),
  \label{eq:dofs}
\end{equation}
giving 3 scalar DOFs $(p,q,r)$ per vertex.  Principal stresses are eigenvalues
of the $2\times 2$ matrix $\bigl[\begin{smallmatrix}p&r\\r&q\end{smallmatrix}\bigr]$
divided by $t$.

%% -------------------------------------------------------------------------
\subsection{Generalised finite-difference (GFDM) operators}

At each vertex $i$ with local frame $(v_1^{(i)},v_2^{(i)},\nn^{(i)})$ we
gather the $k$-ring neighbourhood $\mathcal{N}(i)$ (ring depth $d=3$ by
default, giving $\sim\!30$ neighbours).  Neighbour positions are projected into
the tangent plane as $(\xi_j,\eta_j)$ coordinates.  We then solve a
weighted least-squares quadratic Taylor expansion
\begin{equation}
  f_j - f_i \;\approx\;
    a_i\xi_j + b_i\eta_j
    + \tfrac12 c_i\xi_j^2 + d_i\xi_j\eta_j + \tfrac12 e_i\eta_j^2,
  \label{eq:taylor}
\end{equation}
for the 5 coefficients $(a_i,\dots,e_i)$.  The first-order coefficients give
\emph{stencil weights} that assemble into sparse $n\times n$ matrices:
\begin{equation}
  (G_\xi f)_i \approx \frac{\partial f}{\partial\xi}\Big|_i,
  \qquad
  (G_\eta f)_i \approx \frac{\partial f}{\partial\eta}\Big|_i.
  \label{eq:gfdm-ops}
\end{equation}
Similarly, the second-order coefficients yield $H_{\xi\xi},H_{\xi\eta},H_{\eta\eta}$.

%% -------------------------------------------------------------------------
\subsection{Assembly and solve}

Substituting~\eqref{eq:dofs} into~\eqref{eq:divs-ambient} and
\eqref{eq:equil} yields a sparse linear system $L\,\vv{s}=\vv{b}$ of size
$3n\times 3n$:
\begin{equation}
  L\,\vv{s} + \vv{b} = \vv{0}, \qquad
  \vv{b}^{(i)} = -\Delta p\,\nn^{(i)}, \qquad
  \vv{s}^{(i)} = (p_i,q_i,r_i).
  \label{eq:system}
\end{equation}
Because the system is \emph{rank-deficient} on a closed surface (null modes),
we solve the Tikhonov-regularised normal equations:
\begin{equation}
  (L^\top L + \lambda^2 R^\top R)\,\vv{s} = L^\top\vv{b},
  \label{eq:tikhonov}
\end{equation}
with $R$ a smoothness regulariser (e.g.\ identity or discrete Laplacian) and
$\lambda\approx 0.01$--$0.05$.  After solving, a light Laplacian
post-smoothing ($\sigma\leftarrow(1-\alpha)\sigma+\alpha\bar\sigma_{1\text{-ring}}$,
$\alpha=0.5$, $\sim\!12$ iterations) suppresses the residual icosahedral
discretisation artefacts.

%% =========================================================================
\section{Extracting the principal curvature frame}
\label{sec:pc-frame}

%% -------------------------------------------------------------------------
\subsection{Shape operator from the Weingarten map}

For each vertex $i$ with local frame $(v_1,v_2,\nn)$, the GFDM operators give
the surface gradient of the normal field $\nn$:
\begin{equation}
  \frac{\partial\nn}{\partial\xi}\Big|_i = G_\xi\,\nn, \qquad
  \frac{\partial\nn}{\partial\eta}\Big|_i = G_\eta\,\nn.
  \label{eq:grad-n}
\end{equation}
The components of the second fundamental form in the local frame are then
\begin{align}
  B_{11} &= \frac{\partial\nn}{\partial\xi}\cdot v_1, &
  B_{22} &= \frac{\partial\nn}{\partial\eta}\cdot v_2, &
  B_{12} &= \tfrac12\!\left(\frac{\partial\nn}{\partial\xi}\cdot v_2
                           + \frac{\partial\nn}{\partial\eta}\cdot v_1\right).
  \label{eq:B-components}
\end{align}
(Sign: $+\mathrm{grad}_s\nn$, so $\mathrm{tr}\,B=2H>0$ for an outward normal
under positive pressure.)

%% -------------------------------------------------------------------------
\subsection{Analytic 2$\times$2 diagonalisation}

The shape operator is represented by the symmetric $2\times 2$ matrix
$\mathbf{B}=\bigl[\begin{smallmatrix}B_{11}&B_{12}\\B_{12}&B_{22}\end{smallmatrix}\bigr]$
in the $(v_1,v_2)$ basis.  Its eigenvalues are the principal curvatures:
\begin{equation}
  H_{\rm loc} = \tfrac12(B_{11}+B_{22}), \quad
  d = \sqrt{\tfrac14(B_{11}-B_{22})^2 + B_{12}^2},
  \label{eq:eig-Hloc}
\end{equation}
\begin{equation}
  \kone = H_{\rm loc}+d \;\ge\; \ktwo = H_{\rm loc}-d.
  \label{eq:eig-kappas}
\end{equation}
The eigenvector for $\kone$ makes angle $\theta$ with $v_1$ in the tangent
plane, where
\begin{equation}
  \theta = \tfrac12\arctan\!\frac{2B_{12}}{B_{11}-B_{22}}.
  \label{eq:theta}
\end{equation}
The principal curvature directions in world $\R^3$ are then
\begin{equation}
  \vv{e}_1^{\rm pc} = \cos\theta\;v_1 + \sin\theta\;v_2, \qquad
  \vv{e}_2^{\rm pc} = \nn\times\vv{e}_1^{\rm pc}.
  \label{eq:pc-dirs}
\end{equation}
The resulting frame $(\vv{e}_1^{\rm pc},\vv{e}_2^{\rm pc},\nn)$ is orthonormal
and right-handed by construction.  At umbilic points ($d=0$) the angle
$\theta=0$ by convention; any orthonormal pair in the tangent plane is equally
valid.

%% -------------------------------------------------------------------------
\subsection{Sign ambiguity and global consistency}
\label{sec:sign}

\paragraph{The line-field problem.}
Each principal direction is defined only up to sign: $\vv{e}_1^{\rm pc}$ and
$-\vv{e}_1^{\rm pc}$ are physically equivalent.  The $\arctan2$ formula in
equation~\eqref{eq:theta} picks a sign at each vertex \emph{independently}.
Near a \emph{umbilic point} ($\kone\approx\ktwo$, $d\to0$), the discriminant
vanishes and $\arctan2$ is sensitive to floating-point noise: the sign chosen
at vertex $i$ can oppose that of its neighbour $j$, even though the true
direction varies smoothly.  Numerical experiments on the prolate spheroid
(IcoSphere, $\sim$2600 vertices) show that increasing the neighbourhood depth
$d_{\rm ring}$ from 2 to 5 does not reduce the number of sign-flipped vertices
(34--40 at all depths), confirming that the issue is geometric, not a
fit-accuracy problem.

\paragraph{BFS sign propagation.}
A single breadth-first traversal of the mesh graph (from an arbitrary seed
vertex) enforces a globally consistent orientation.  Starting from vertex 0
with its computed $\vv{e}_1^{\rm pc}$, every unvisited neighbour $j$ of a
visited vertex $i$ is checked:
\begin{equation}
  \text{if } \vv{e}_1^{\rm pc}(i)\cdot\vv{e}_1^{\rm pc}(j) < 0,\quad
  \text{then } \vv{e}_1^{\rm pc}(j)\leftarrow -\vv{e}_1^{\rm pc}(j),\;
               \vv{e}_2^{\rm pc}(j)\leftarrow -\vv{e}_2^{\rm pc}(j).
  \label{eq:bfs}
\end{equation}
Flipping both $\vv{e}_1^{\rm pc}$ and $\vv{e}_2^{\rm pc}$ simultaneously
preserves the right-handedness of the frame:
$(-\vv{e}_1^{\rm pc})\times(-\vv{e}_2^{\rm pc})=\vv{e}_1^{\rm pc}\times\vv{e}_2^{\rm pc}=\nn$.
On the prolate spheroid, BFS reduces the number of sign-flipped 1-ring pairs
from 34 to 2.

\paragraph{Unavoidable topological singularities.}
The 2 residual flipped vertices correspond to the two \emph{umbilic points} at
the tips of the spheroid, where $\kone=\ktwo$ exactly.  These are genuine
topological singularities of the principal curvature line field: the
eigenvectors of $\mathbf{B}$ wind around such a point and no globally
consistent orientation exists in its neighbourhood.  This is not a numerical
artefact.

The Poincaré--Hopf theorem, applied to the principal curvature line field,
requires that the sum of singularity indices equals the Euler characteristic
$\chi(\Gamma)$ of the surface.  For a closed genus-0 surface
$\chi=2$; each regular umbilic point contributes index $+1$ to this sum.  The
prolate spheroid has exactly two umbilic points (both poles), each of index
$+1$, saturating the bound.  On an arbitrary smooth closed surface of genus 0
the Carath\'eodory conjecture asserts that there are always at least two
umbilic points; hence \textbf{at least 2 sign-propagation singularities are
unavoidable on any closed genus-0 mesh}, regardless of resolution or
neighbourhood size.

%% =========================================================================
\section{Summary: the per-vertex output}
\label{sec:output}

For each vertex the pipeline produces:

\begin{center}
\begin{tabular}{lll}
\toprule
\textbf{Quantity} & \textbf{Symbol} & \textbf{Dimensions} \\
\midrule
Principal curvatures    & $\kone,\ktwo$                & $\mathrm{m}^{-1}$ \\
Principal radii         & $R_1=1/\kone,\;R_2=1/\ktwo$ & $\mathrm{m}$       \\
Mean / Gaussian         & $H,K$                        & $\mathrm{m}^{-1}$, $\mathrm{m}^{-2}$ \\
Principal direction 1   & $\vv{e}_1^{\rm pc}\in\R^3$  & ---                \\
Principal direction 2   & $\vv{e}_2^{\rm pc}\in\R^3$  & ---                \\
Outward unit normal     & $\nn\in\R^3$                 & ---                \\
\midrule
Stress resultant (after solve)  & $N_{11},N_{22},N_{12}$ & $\mathrm{N/m}$ \\
Principal stresses              & $\sigma_1,\sigma_2$     & $\mathrm{Pa}$  \\
\bottomrule
\end{tabular}
\end{center}

\noindent The curvature frame is computed by \texttt{surface\_curvature\_frame.py}
(\texttt{compute\_curvature\_frame}); the stress solve is in
\texttt{membrane\_stress\_fd.py} (\texttt{solve\_membrane}).

%% =========================================================================
\section{Validation cases}
\label{sec:validation}

\paragraph{Sphere, $R=1$ (totally umbilic — worst case).}
All principal curvatures equal $1/R=1$; $H=1$, $K=1$ everywhere.  Because
$d\equiv0$ at every vertex, every point is umbilic and principal directions are
completely undefined: $\mathbf{B}=(1/R)\mathbf{I}_2$ so the $\arctan2$ formula
receives $(0,\,0)$ and the sign at each vertex is set by floating-point noise in
the GFDM estimate of $B_{12}$.  Frame orthonormality is still preserved to
machine precision ($|\vv{e}_1\cdot\vv{e}_2|\le10^{-16}$), but the
\emph{directions} are numerically arbitrary.

After BFS propagation the result is a smooth tangent vector field on $S^2$.
By the \textbf{hairy ball theorem} (a corollary of Poincaré--Hopf for $S^2$)
no continuous non-vanishing tangent vector field can exist on the sphere; the
BFS field must therefore develop singularities.  On an IcoSphere subdiv-4
($n=2562$ vertices), 314 vertices (12\% of the mesh) have more than 40\% of
their 1-ring neighbours pointing in the opposite direction after BFS —
forming broad singular \emph{patches} rather than isolated points, because the
entire surface is critical.  The curvature scalars ($\kone,\ktwo,H,K$) are
unaffected; only the directional output (\textbf{e}$_1$, \textbf{e}$_2$) is
meaningless on a totally umbilic surface.

\paragraph{Prolate spheroid, $a=2$, $b=1$ (long axis $x$).}
Away from the poles the surface is non-umbilic ($d>0$) and the principal
directions are well-defined.  At the equator ($|x|<0.1a$, $\sim\!1036$
vertices on IcoSphere subdiv-4):

\begin{center}
\begin{tabular}{lrrr}
\toprule
Quantity & Computed (mean) & Std & Analytic \\
\midrule
$\kone$ (hoop)       & 1.004 & 0.028 & $1/b=1.000$ \\
$\ktwo$ (meridional) & 0.248 & 0.007 & $b/a^2=0.250$ \\
\bottomrule
\end{tabular}
\end{center}

Frame orthonormality: $|\vv{e}_1\cdot\vv{e}_2|\le2\times10^{-16}$ (machine
precision).  Membrane stresses: $\sigma_{\rm hoop}=\Delta p r_2(1-r_2/2r_1)/t$,
$\sigma_{\rm merid}=\Delta p r_2/2t$.

\paragraph{Capsule, $R=1$, $H=2$ (cylinder + hemispherical end-caps).}

A capsule is a cylinder of radius $R$ and half-height $H$ capped with two
hemispheres.  It has three geometrically distinct regions:
\begin{itemize}[nosep]
  \item \textbf{Cylindrical body} ($|z|<H$): $\kone=1/R$ (hoop), $\ktwo=0$ (axial).
        Non-umbilic; $d=1/(2R)$; principal directions well-defined.
  \item \textbf{Spherical caps} ($|z|>H+R\,\varepsilon$): $\kone=\ktwo=1/R$.
        Totally umbilic; $d=0$; same sign-propagation pathology as the sphere.
  \item \textbf{Junction} ($|z|\approx H$): $d$ ramps smoothly between the two values.
\end{itemize}
The parametric capsule mesh (\texttt{make\_capsule} in \texttt{show\_capsule.py},
$n_\theta=40$, $n_\varphi=14$, matched ring spacing) gives $n=2522$ vertices.

\begin{center}
\begin{tabular}{lrrrrr}
\toprule
Region & $n$ & $\bar\kone$ & $\sigma(\kone)$ & $\bar\ktwo$ & $\sigma(\ktwo)$ \\
\midrule
Cylinder body ($n_{\rm cyl}=1400$) & analytic & \multicolumn{2}{c}{$\kone=1.000$} & \multicolumn{2}{c}{$\ktwo=0$} \\
 & computed & 1.000 & 0.000 & 0.034 & 0.083 \\[2pt]
Spherical caps ($n_{\rm cap}=1042$) & analytic & \multicolumn{2}{c}{$\kone=1.000$} & \multicolumn{2}{c}{$\ktwo=1.000$} \\
 & computed & 1.037 & 0.047 & 0.895 & 0.076 \\
\bottomrule
\end{tabular}
\end{center}

The cylinder $\kone$ is exact (GFDM error $<0.001$); $\ktwo=0.034$ (3.4\%)
reflects the stencil mixing axial and hoop contributions near the coarse
junction seam.  Cap errors ($\sim$4--10\%) are highest at cap vertices whose
stencils span the junction.

Sign-consistency after BFS: \textbf{3} inconsistent vertices out of 2522
(0.1\%), located at the two umbilic poles and one nearby junction vertex.

\paragraph{Sign-consistency comparison ($n\approx2550$, after BFS).}

\begin{center}
\begin{tabular}{lrrr}
\toprule
 & \textbf{Sphere} & \textbf{Spheroid} & \textbf{Capsule} \\
 & (all umbilic)   & (2 umbilic poles) & (2 umbilic poles + cylinder) \\
\midrule
Inconsistent vertices ($>$40\% flipped) & 314\;(12\%) & 2\;($<$0.1\%) & 3\;(0.1\%) \\
Max flip fraction                        & 0.83 & 0.67 & 0.67 \\
Mean flip fraction                       & 0.12 & 0.018 & 0.021 \\
Depth sweep (2--5) helps?               & no & no & no \\
\bottomrule
\end{tabular}
\end{center}

The capsule behaves like the spheroid: BFS reduces sign flips to the 2--3
unavoidable singularities at the two umbilic cap poles.  The cylinder body is
entirely clean.  For real biological meshes (non-umbilic almost everywhere) the
spheroid/capsule result is the relevant model.  The sphere is the degenerate
worst case.

\end{document}
